\documentclass[11pt]{article}
\usepackage[margin=1in]{geometry}
\usepackage{amsmath,amssymb,amsthm}
\usepackage{graphicx}
\usepackage[hidelinks]{hyperref}

\newtheorem{theorem}{Theorem}
\newtheorem{lemma}{Lemma}
\newtheorem{corollary}{Corollary}

\newcommand{\A}{\mathcal A}
\newcommand{\B}{\mathcal B}
\newcommand{\Hh}{\mathcal H}
\newcommand{\K}{\mathcal K}
\newcommand{\M}{\mathcal M}
\newcommand{\N}{\mathbb N}
\newcommand{\Oc}{\mathcal O}
\newcommand{\Q}{\mathcal Q}
\newcommand{\C}{\mathbb C}
\newcommand{\id}{\mathrm{id}}
\newcommand{\UCP}{\mathrm{UCP}}
\newcommand{\Cstar}{C^*}

\title{A Counterexample to Compact Refinement for the Approximate Unique Extension Property}
\author{fa\_banach\_001, agent\_lane\_06}
\date{2026-07-03}

\begin{document}
\maketitle

\paragraph{Status.}
Candidate counterexample, likely valid.  This packet answers negatively the
compact-refinement question raised after Theorem 3.6 of Thompson's paper
\cite{ThompsonAUEP}.

\section{Source and Question}

The source paper studies the approximate unique extension property (AUEP) for
unital completely positive extensions of representations of operator spaces.
After proving that approximate unitary equivalence to a representation with the
unique extension property implies AUEP, Thompson compares the result with
Voiculescu's theorem and asks whether, under separability assumptions, the
approximating errors
\[
   \pi(t)-u_\beta^*\psi(t)u_\beta
\]
can always be chosen compact for every $t$.

The relevant source passage is on pages 11--12 of \cite{ThompsonAUEP}.

\begin{center}
\includegraphics[width=\textwidth]{figures/open_problem_crop.png}
\end{center}

\section{Idea of the Proof}

The compactness refinement is too strong because compact perturbation is visible
in the Calkin algebra.  If a unitary $u$ made
$\pi(t)-u^*\psi(t)u$ compact for every $t$, then the Calkin image of $\psi$
would be unitarily conjugate to the Calkin image of the representation $\pi$.
Thus the Calkin image of $\psi$ would have to be multiplicative.

Now start with an AUEP representation that is not UEP.  It has a nontrivial UCP
extension, hence a nonmultiplicative one.  Countably amplify both maps.  The
multiplicative defect becomes an infinite diagonal repeat of a nonzero operator,
so it is not compact.  This contradicts the Calkin multiplicativity forced by
compact errors.  Thompson's Cuntz-algebra example supplies the required
separable AUEP-but-not-UEP representation.

\section{A General Amplification Obstruction}

Let $\Q(\Hh)=\B(\Hh)/\K(\Hh)$ denote the Calkin algebra and let
$q:\B(\Hh)\to\Q(\Hh)$ be the quotient map.

\begin{lemma}[Calkin obstruction]\label{lem:calkin}
Let $A$ be a unital $C^*$-algebra, let $\pi:A\to\B(\Hh)$ be a
$*$-representation, and let $\psi:A\to\B(\Hh)$ be a UCP map.  If there is a
unitary $u\in\B(\Hh)$ such that
\[
   \pi(t)-u^*\psi(t)u\in \K(\Hh) \qquad (t\in A),
\]
then $q\circ\psi:A\to\Q(\Hh)$ is a $*$-homomorphism.
\end{lemma}

\begin{proof}
In the Calkin algebra,
\[
   q(\pi(t))=q(u)^*q(\psi(t))q(u) \qquad (t\in A).
\]
Hence
\[
   q(\psi(t))=q(u)q(\pi(t))q(u)^* \qquad (t\in A).
\]
The right-hand side is a unitary conjugate of the $*$-homomorphism $q\circ\pi$,
so $q\circ\psi$ is a $*$-homomorphism.
\end{proof}

\begin{theorem}[Amplified non-UEP obstruction]\label{thm:general}
Let $\M$ be a separable unital operator space and let
$\pi_0:\Cstar(\M)\to\B(\Hh_0)$ be a representation on a separable Hilbert space.
Assume that $\pi_0$ has AUEP with respect to $\M$ but does not have UEP.
Then there are a separable Hilbert space $\Hh$, a representation
$\Pi:\Cstar(\M)\to\B(\Hh)$ with AUEP, and a UCP map
$\Psi:\Cstar(\M)\to\B(\Hh)$ such that $\Psi|_\M=\Pi|_\M$, but no unitary
$u\in\B(\Hh)$ satisfies
\[
   \Pi(t)-u^*\Psi(t)u\in\K(\Hh) \qquad (t\in \Cstar(\M)).
\]
Consequently no AUEP-approximating net or sequence for this pair can have all
of its pointwise errors compact.
\end{theorem}

\begin{proof}
Because $\pi_0$ is not UEP, there exists a UCP map
$\psi_0:\Cstar(\M)\to\B(\Hh_0)$ such that
$\psi_0|_\M=\pi_0|_\M$ and $\psi_0\ne\pi_0$.  The map $\psi_0$ cannot be
multiplicative: if it were a $*$-homomorphism, then its agreement with $\pi_0$
on the generating space $\M$ would force $\psi_0=\pi_0$ on $\Cstar(\M)$.
Thus there exist $a,b\in\Cstar(\M)$ such that
\[
   D=\psi_0(ab)-\psi_0(a)\psi_0(b)\ne 0.
\]

Let $\Hh=\bigoplus_{n=1}^\infty \Hh_0$, and define the countable amplifications
\[
   \Pi=\bigoplus_{n=1}^\infty \pi_0,\qquad
   \Psi=\bigoplus_{n=1}^\infty \psi_0.
\]
The Hilbert space $\Hh$ is separable.  The representation $\Pi$ has AUEP by
Thompson's direct-sum permanence theorem for AUEP
\cite[Proposition 3.7]{ThompsonAUEP}, and $\Psi|_\M=\Pi|_\M$.

The multiplicative defect of $\Psi$ at $a,b$ is
\[
   \Delta=\Psi(ab)-\Psi(a)\Psi(b)=\bigoplus_{n=1}^\infty D.
\]
Since $D\ne0$, choose $\xi\in\Hh_0$ with $\|D\xi\|>0$.  The unit vectors placing
$\xi/\|\xi\|$ in the $n$th summand show that $\Delta$ maps a bounded sequence
to a sequence with no norm-convergent subsequence.  Therefore $\Delta$ is not
compact.

If a unitary $u$ made all errors $\Pi(t)-u^*\Psi(t)u$ compact, then
Lemma \ref{lem:calkin} would imply that $q\circ\Psi$ is a $*$-homomorphism.
In particular,
\[
   q(\Delta)
   =q(\Psi(ab)-\Psi(a)\Psi(b))
   =0.
\]
That says exactly that $\Delta$ is compact, contradicting the previous
paragraph.  Hence no such unitary exists.
\end{proof}

\section{Concrete Counterexample from the Source Paper}

\begin{corollary}\label{cor:cuntz}
The compact-refinement question raised in \cite{ThompsonAUEP} has a negative
answer even when both the operator space and the Hilbert space are separable.
\end{corollary}

\begin{proof}
Let $\Oc_\infty$ be the infinite Cuntz algebra with generators $(V_n)_{n\in\N}$,
and let $\A$ be the norm-closed unital operator algebra generated by these
generators.  Thompson recalls that $\Cstar(\A)\cong\Oc_\infty$, that all
irreducible representations of $\Oc_\infty$ are approximately unitarily
equivalent because $\Oc_\infty$ is simple, and that there is an irreducible
representation $\rho$ which is not a boundary representation for $\A$
\cite[Example 2]{ThompsonAUEP}.  Since $\rho$ is not a boundary representation,
it does not have UEP with respect to $\A$.

The same example supplies boundary representations, and all irreducible
representations are approximately unitarily equivalent.  Therefore $\rho$ is
approximately unitarily equivalent to a representation with UEP.  By Thompson's
Theorem 3.6, $\rho$ has AUEP.  Because $\Oc_\infty$ is separable and $\rho$ is
irreducible, $\rho$ acts on a separable Hilbert space: any nonzero vector is
cyclic and the orbit of a countable dense subset of $\Oc_\infty$ is dense.

Apply Theorem \ref{thm:general} to $\M=\A$ and $\pi_0=\rho$.  The resulting
countable amplification is a separable AUEP representation with a UCP extension
for which compact pointwise errors are impossible.  This is precisely a negative
instance of the compact-refinement property under separability assumptions.
\end{proof}

\section{Verification Notes and Limitations}

The proof is purely structural; no numerical computation is used.  The main
check is that compact pointwise errors imply multiplicativity in the Calkin
quotient, whereas countable amplification turns any nonzero multiplicative
defect into a noncompact diagonal defect.

The result answers the general compact-refinement question negatively.  It does
not rule out positive refinements under additional hypotheses, such as UEP
itself, postliminal envelope assumptions, or restrictions on the Calkin image of
the chosen UCP extension.

\section{Novelty Check}

The run's lightweight indexes were searched for the arXiv id, the exact title,
``approximate unique extension property'', ``compact refinement'', ``Calkin'',
and the Voiculescu compact-error phrase.  No previous packet for this exact
target was found.  A web search on 2026-07-03 for exact close variants of
``approximate unique extension property'' with ``compact'', ``Calkin'', and
``Voiculescu'' found the source arXiv paper but no separate resolution.
Novelty confidence is therefore bounded-local-plus-web rather than exhaustive.

\section{Human Review Recommendation}

Review as a candidate counterexample.  The most important points to check are:
(1) Thompson's Example 2 really supplies a separable AUEP representation that is
not UEP; (2) the non-UEP extension is necessarily nonmultiplicative; and (3) the
Calkin-quotient argument forbids even one compact-error conjugating unitary.

\begin{thebibliography}{9}

\bibitem{ThompsonAUEP}
Ian Thompson.
\newblock \emph{An approximate unique extension property for completely positive maps}.
\newblock arXiv:2301.08285v1, 2023.

\end{thebibliography}

\end{document}
